% TOP_PROBLEM: {"id": 3134, "title": "Fractional Hadwiger", "category_id": 3, "category": {"id": 3, "name": "graph_theory", "display_name": "Graph Theory", "description": "Problems involving graphs, networks, and their properties.", "slug": "graph-theory", "order_index": 3, "created_at": "2026-07-31T15:26:25.671Z"}, "status": "open", "problem_number": "OPG-59908", "set_id": 12}
% FIRST_POSTED: 2026-10-01
% ENTRY_KIND: statement_only
% UnsolvedMath ID 3134; source catalogue code OPG-59908
% !TeX program = lualatex
% Definition and sources only; this file is not an LLM solution attempt.
\documentclass[11pt,a4paper]{article}
\usepackage[margin=25mm]{geometry}
\usepackage{fontspec}
\setmainfont{lmroman10-regular.otf}[BoldFont=lmroman10-bold.otf,
 ItalicFont=lmroman10-italic.otf,BoldItalicFont=lmroman10-bolditalic.otf]
\usepackage{amsmath,amssymb,mathtools}
\usepackage{unicode-math}
\setmathfont{latinmodern-math.otf}
\AtBeginDocument{\let\setminus\smallsetminus}
\usepackage{xurl,hyperref}
\hypersetup{colorlinks=true,urlcolor=blue}
\setcounter{secnumdepth}{0}
\setlength{\parindent}{0pt}
\setlength{\parskip}{5pt}
\setlength{\emergencystretch}{3em}
\begin{document}
\section*{Fractional Hadwiger}\label{3134}
{\small UnsolvedMath ID 3134 · OPG-59908 · Graph Theory · OpenGarden}
\subsection{Definitions and mathematical statement}
Let $G$ be a finite nonempty simple graph. Its chromatic number $\chi(G)$ is the least number of independent vertex sets partitioning $V(G)$. Its Hadwiger number $h(G)$ is the largest $t$ for which there are $t$ pairwise disjoint nonempty connected vertex sets, every pair joined by an edge. Such sets are branch sets of a complete minor.

Define the fractional chromatic number $\chi_f(G)$ by assigning nonnegative real weights $x_I$ to independent sets and minimizing $\sum_Ix_I$ subject to $\sum_{I\ni v}x_I\ge1$ for every vertex $v$. For the fractional Hadwiger number, a bramble is a family $\mathcal B$ of nonempty connected vertex sets such that every two touch: they either intersect or an edge joins them. Assign nonnegative real weights $y_B$ satisfying $\sum_{B\ni v}y_B\le1$ at each vertex. Let $h_f(G)$ be the maximum of $\sum_{B\in\mathcal B}y_B$ over all such brambles and weights. Since $G$ is finite, these optimization problems involve finitely many possible vertex sets.

The source asks for the validity, for every such $G$, of each of the following inequalities:
\[
 \chi_f(G)\le h(G),\qquad
 \chi(G)\le h_f(G),\qquad
 \chi_f(G)\le h_f(G).
\]
These are separate assertions. In the fractional minor parameter the touching condition is imposed on the entire positive-weight support, rather than merely on some pairs or in expectation; omitting it would define a different quantity.
\subsection{Short English statement}
Compare ordinary and fractional coloring requirements with ordinary and fractional complete-minor parameters in each of three proposed inequalities.
\subsection{Sources}
\begin{itemize}
\item[S1] ULAM.AI and the UnsolvedMath Contributors, supplied problems.json, record 3134 (OPG-59908), OpenGarden. Catalogue identity and source transcription; CC BY 4.0.
\url{https://huggingface.co/datasets/ulamai/UnsolvedMath}.
\item[S2] Open Problem Garden, Fractional Hadwiger. Original problem and discussion.
\url{http://www.openproblemgarden.org/op/fractional_hadwiger}.
\end{itemize}
\end{document}
